\providecommand{\P}{\mathbb{P}}
\providecommand{\C}{\mathbb{C}}
\providecommand{\Z}{\mathbb{Z}}
\providecommand{\cO}{\mathcal{O}}
\providecommand{\cH}{\mathcal{H}}
\providecommand{\fgl}{\mathfrak{gl}}
\providecommand{\CoHA}{\mathrm{CoHA}}
\providecommand{\Tot}{\mathrm{Tot}}
\providecommand{\Crit}{\mathrm{Crit}}
\providecommand{\tr}{\mathrm{tr}}
\providecommand{\Spec}{\mathrm{Spec}}

\section*{Local $\P^2$ from three affine charts and the McKay positive half}
\label{sec:local-p2-three-charts}

Let
\[
 X = K_{\P^2}=\Tot(\cO_{\P^2}(-3)).
\]
This is a non-compact toric CY$_3$ and a local surface geometry: the
total space of the canonical line over the surface $\P^2$. The output of
$\Phi_3$ belongs to the $E_1$ chiral sector after specialisation; an
$E_2$ structure belongs to the centre, double, module, or representation
layer, with extra pairing and completion data. The symbol $Y^+(X)$ below
denotes the Hall positive half attached to a fixed McKay heart.

\subsection*{Crepant fan and affine cover}

Let $N=\Z^3$ and write
\[
 v_0=(0,0,1),\qquad
 v_1=(1,0,1),\qquad
 v_2=(0,1,1),\qquad
 v_3=(-1,-1,1).
\]
The determinant of the singular cone $\langle v_1,v_2,v_3\rangle$ is
$3$, and its quotient type is $\C^3/(\Z/3)$ for the diagonal action of
weights $(1,1,1)$.
The crepant fan of $K_{\P^2}$ inserts the interior ray $v_0$ and has
maximal cones
\[
 \sigma_0=\langle v_0,v_1,v_2\rangle,\qquad
 \sigma_1=\langle v_0,v_2,v_3\rangle,\qquad
 \sigma_2=\langle v_0,v_3,v_1\rangle .
\]
Each determinant is $\pm1$, hence
\[
 U_i=\Spec \C[\sigma_i^\vee\cap M]\simeq \C^3.
\]
The pairwise intersections are the affine toric charts of the common
walls,
\[
 U_0\cap U_1\simeq U_{\langle v_0,v_2\rangle}\simeq \C^2\times\C^*,
 \quad
 U_1\cap U_2\simeq U_{\langle v_0,v_3\rangle}\simeq \C^2\times\C^*,
 \quad
 U_2\cap U_0\simeq U_{\langle v_0,v_1\rangle}\simeq \C^2\times\C^*.
\]
The triple intersection is
\[
 U_0\cap U_1\cap U_2\simeq U_{\langle v_0\rangle}
 \simeq \C\times(\C^*)^2.
\]
All rays have height one with respect to the character
$m_\Omega=(0,0,1)\in M$, so the toric canonical divisor is principal.
This character gives the global holomorphic volume form; in the affine
coordinates of every $U_i\simeq\C^3$ it is the standard form multiplied
by a unit on overlaps.

\subsection*{McKay heart and NCCR}

Let $G=\Z/3$ act on $\C^3$ by
\[
 \zeta\cdot(x,y,z)=(\zeta x,\zeta y,\zeta z).
\]
The skew group algebra $\C[x,y,z]\#G$ is presented by the McKay quiver
with vertices $\Z/3$ and arrows
\[
 a_i^{(x)},a_i^{(y)},a_i^{(z)}\colon i\longrightarrow i+1,
 \qquad i\in\Z/3.
\]
Its cubic potential is
\[
 W_{K_{\P^2}}
 =\sum_{i\in\Z/3}\tr\!\left(
 a_i^{(x)}a_{i+1}^{(y)}a_{i+2}^{(z)}
 -a_i^{(x)}a_{i+1}^{(z)}a_{i+2}^{(y)}
 \right).
\]
The cyclic derivatives give the commutative square relations
\[
 a_i^{(u)}a_{i+1}^{(v)}=a_i^{(v)}a_{i+1}^{(u)},
 \qquad u,v\in\{x,y,z\},\quad i\in\Z/3.
\]
Thus
\[
 \mathrm{Jac}(Q_{K_{\P^2}},W_{K_{\P^2}})
 \cong \C[x,y,z]\#(\Z/3).
\]
By the derived McKay equivalence of Bridgeland, King, and Reid
\cite{BridgelandKingReid2001}, the $G$-Hilbert crepant resolution
$K_{\P^2}\to\C^3/G$ has derived category equivalent to the derived
category of modules over this Jacobi algebra. Kapranov's exceptional
collection $\{\cO,\cO(1),\cO(2)\}$ on $\P^2$ \cite{Kapranov1988} lifts
to the tilting bundle
\[
 T=\pi^*(\cO\oplus\cO(1)\oplus\cO(2))
\]
on $K_{\P^2}$. The algebra $\mathrm{End}_{K_{\P^2}}(T)$ is the
Morita-equivalent NCCR obtained from the same McKay relations.

\subsection*{Local positive halves}

Each affine chart $U_i\simeq\C^3$ carries the Jordan quiver with cubic
potential. Its critical CoHA is the Schiffmann and Vasserot positive
half
\[
 \CoHA(U_i)=\cH(Q_{\C^3},W_{\C^3})
 \cong Y^+(\widehat{\fgl}_1)
\]
as an associative Hall algebra \cite{SV2013}. On an overlap
$\C^2\times\C^*$ one coordinate has been inverted. The shuffle
description of $Y^+(\widehat{\fgl}_1)$ is therefore localised by the
corresponding monomial, and the wall transition is the
Kontsevich and Soibelman automorphism between the adjacent toric
chambers. These maps preserve the positive Hall half.

The global algebra used for local $\P^2$ is the NCCR critical CoHA
\[
 \cH_{K_{\P^2}}
 =
 \bigoplus_{\mathbf d\in\Z_{\geq0}^3}
 H^\bullet_{G_{\mathbf d}}\!
 \left(\Crit(\tr W_{\mathbf d}),\varphi_{\tr W_{\mathbf d}}\right),
 \qquad
 G_{\mathbf d}=\prod_{i\in\Z/3}\mathrm{GL}_{d_i}.
\]
We write
\[
 Y^+(K_{\P^2}) := \cH_{K_{\P^2}}.
\]
This definition stays inside the positive half. It is a quiver-with-
potential construction. A \v{C}ech equaliser presentation requires
orientation data, overlap Morita kernels, the cocycle condition on the
triple intersection, and a recognition theorem identifying the resulting
Hall algebra with the NCCR CoHA. The Drinfeld double
$D(Y^+(K_{\P^2}))$ adjoins Cartan and negative generators only after a
Hopf pairing, PBW theorem, completion, and centre compatibility are
supplied.

\subsection*{Boundary of the RSYZ theorem}

Rapcak, Soibelman, Yang, and Zhao identify critical CoHAs with positive
halves of affine super Yangians for the toric CY$_3$ geometries in their
stated range \cite{RSYZ2020}. The Schiffmann--Vasserot theorem supplies
the $\C^3$ chart normal form, and the RSYZ comparison supplies the
positive-half dictionary where its hypotheses hold. Local $\P^2$
contains the compact divisor $\P^2$, equivalently the toric diagram has
the interior ray $v_0$. For this geometry the constructed global object
in this note is the McKay NCCR CoHA above; a named affine super Yangian
presentation is an additional descent-and-recognition theorem.

Consequently the symbol $Y^+(K_{\P^2})$ records a positive Hall algebra.
An affine super Yangian presentation, a full quantum group, or a
$\mathcal{W}$-algebra realisation requires the additional double and
centre data named above.

\subsection*{\texorpdfstring{$\kappa_{\mathrm{ch}}$}{kch} of local
\texorpdfstring{$\P^2$}{P2}}

The chiral invariant of the local surface is
\[
 \kappa_{\mathrm{ch}}(K_{\P^2})
 =\frac{1}{2}\chi_{\mathrm{top}}(\P^2)
 =\frac{3}{2}.
\]
This is the local $\P^2$ row of
Theorem~\ref{thm:local-p2-shadow} in the Vol III
\(\kappa_\bullet\)-stratification chapter.

The remaining three invariants are distinct from this value. The
base-normalised categorical witness is
\[
 \chi(\cO_{\P^2})=h^0(\cO_{\P^2})-h^1(\cO_{\P^2})+h^2(\cO_{\P^2})
 =1.
\]
Since $K_{\P^2}$ is non-compact, a smooth-proper
$\kappa_{\mathrm{cat}}(K_{\P^2})$ requires an extra convention; the
displayed number is a base witness. The Borcherds invariant
$\kappa_{\mathrm{BKM}}(K_{\P^2})$ is undefined until a Borcherds
denominator for this geometry is constructed. With the local-surface
fibre convention one may record
\[
 \kappa_{\mathrm{fiber}}(K_{\P^2})=\mathrm{rk}\,H_2(\P^2,\Z)=1.
\]
Thus the local-surface convention records the separated tuple
\[
 \bigl(\kappa_{\mathrm{ch}},\kappa_{\mathrm{cat}},
 \kappa_{\mathrm{BKM}},\kappa_{\mathrm{fiber}}\bigr)
 =
 \left(\frac32,\ 1,\ \text{undefined},\ 1\right),
\]
with the second entry understood as the base-normalised witness.

Three first-principles checks give the same value.
\[
 \chi_{\mathrm{top}}(\P^2)=1+1+1=3,
 \qquad
 \frac12\chi_{\mathrm{top}}(\P^2)=\frac32.
\]
The toric cover gives
\[
 \chi_c(X)
 =\sum_i\chi_c(U_i)
 -\sum_{i<j}\chi_c(U_i\cap U_j)
 +\chi_c(U_0\cap U_1\cap U_2)
 =3-0+0=3,
\]
because $\chi_c(\C^3)=1$ and every overlap contains a factor $\C^*$.
The McKay heart has three degree-one simple sectors, one for each
character of $\Z/3$. The critical locus in each simple sector is a point
with one-dimensional vanishing-cycle contribution, so the primitive BPS
index is $\Omega_1=3$ and the chiral normalisation gives
$\Omega_1/2=3/2$.

The resolved conifold has a two-vertex McKay/Klebanov--Witten heart and
lies outside local-surface geometry:
$\Tot(\cO_{\P^1}(-1)\oplus\cO_{\P^1}(-1))$ is a rank-two bundle over a
curve, whereas a local surface has the form $\Tot(\omega_S)$ over a
surface. Direct McKay computation gives
\[
 \kappa_{\mathrm{ch}}(X_{\mathrm{con}})=1.
\]
Local $\P^2$ has the three-vertex $\Z/3$ McKay heart, compact divisor
$\P^2$, and $\kappa_{\mathrm{ch}}(K_{\P^2})=3/2$.

\subsection*{Positive-half statement}

The three affine charts of $K_{\P^2}$ each carry the
$\C^3$ positive half $Y^+(\widehat{\fgl}_1)$. Their overlap maps are
localised shuffle wall-crossing maps. The NCCR positive half used for
the local surface is
\[
 Y^+(K_{\P^2})
 =
 \CoHA\!\left(\mathrm{Jac}(Q_{K_{\P^2}},W_{K_{\P^2}})\right),
\]
with the Jacobi algebra fixed by the Bridgeland, King, and Reid McKay
equivalence and by the Kapranov tilting bundle. The chart character is a
local test; it becomes a global algebra presentation only after descent,
orientation, and recognition data have been supplied. The numerical
invariant attached to this local surface is
$\kappa_{\mathrm{ch}}(K_{\P^2})=3/2$.
